\uuid{FxzP}
\titre{Équation droite/plan}
\auteur{bodin ; hanani ; mzari}
\organisation{exo7}
\licence{CC BY-NC-SA 4.0}
\datecreate{2020-01-20}
\niveau{L1}
\module{Géométrie}
\chapitre{Géométrie de l'espace}
\sousChapitre{Plans}
\difficulte{3}
\theme{Plans -- Droites, équation droite, plan}
\exo7qcmid{lille-1.geometrie-espace.018}

\contenu{
\texte{
Pour ces questions, l'espace est muni d'un repère orthonormé direct  $(O, \vec{i}, \vec{j}, \vec{k})$.
}

\begin{question}
\qclassification{141.03,141.04}
\qtags{motcle=équation droite/plan}
Soit $(P_n), n\in \Nn$, la famille de plans d'équations : $n^2x+(2n-1)y+nz=3$. On note $E$ l'intersection de ces plans, c'est-à-dire $E= \{M(x,y,z) \in \Rr^3; \, M\in P_n, \forall n\in \Nn \}$. Quelles sont les assertions vraies ?
\begin{answers}

       \bad {$E=\emptyset$}

      \bad {$E$ est un plan d'équation $x+y+z=3$.}

     \bad {$E$ est une droite d'équation $ \left\{\begin{array}{ccl}x+y+z&=&3\\y&=&-3 \end{array}\right.$.}

     \good {$E$ est le point de coordonnées $(0,-3,6)$.}

\end{answers}
\begin{explanations}
Soit $M(x,y,z) \in E$, alors   $n^2x+(2n-1)y+nz=3, \, \forall n \in \Nn  \Leftrightarrow  xn^2+ (2y+z)n-y-3=0, \, \forall n \in \Nn  \Leftrightarrow  x=0, 2y+z=0$ et $y+3=0$.
\end{explanations}
\end{question}
}
